\documentclass[11pt,a4paper]{article}

\usepackage[utf8]{inputenc}
\usepackage[T1]{fontenc}
\usepackage[margin=1in]{geometry}
\usepackage{amsmath,amssymb,amsthm,mathtools}
\usepackage{hyperref}
\usepackage{enumitem}
\usepackage{microtype}
\usepackage{booktabs}

\title{Terminal Rigidity Witnesses and the Einstein--Rosen Bridge:\\
A Zero-Capacity Invariant for Physical Inference}
\author{Inacio F. Vasquez\\Independent Researcher}
\date{January 2026}

\newtheorem{theorem}{Theorem}
\newtheorem{lemma}{Lemma}
\newtheorem{definition}{Definition}
\newtheorem{corollary}{Corollary}
<<<<<<< HEAD

\newcommand{\Capacity}{\mathrm{Cap}}
\newcommand{\TC}{\mathrm{TC}}
\newcommand{\SIGC}{\mathbf{SIGC}}
=======
\newtheorem{axiom}{Axiom}

\newcommand{\Capacity}{\mathrm{Cap}}
\newcommand{\TC}{\mathrm{TC}}
\newcommand{\ED}{\mathrm{ED}}
\newcommand{\SIGC}{\mathbf{SIGC}}
\newcommand{\PhysSys}{\mathbf{PhysSys}}
>>>>>>> f178752 (docs: restore canonical ERB LaTeX source)

\begin{document}
\maketitle

\begin{abstract}
<<<<<<< HEAD
We formalize the notion of a \emph{Terminal Rigidity Witness}: a physical system whose global
structure is nontrivial while its operational information capacity is exactly zero.
We show that the classical Einstein--Rosen bridge provides a canonical example.
\end{abstract}

\section{Transcript Capacity}

For an interface producing transcript $Y_{1:T}$ from hidden state $X$, define
\[
\TC(T) := \sup I(X;Y_{1:T}).
\]

\section{Terminal Rigidity}

A system $S$ is a terminal rigidity witness if
\[
\Capacity(\mathcal C_S) = 0
\]
while $S$ has nontrivial global structure.

\section{Einstein--Rosen Bridge}

Topological censorship implies no causal curve connects the two asymptotic regions of the
maximal Schwarzschild extension. Hence the induced channel has zero capacity.

\begin{theorem}
The Einstein--Rosen bridge is a terminal rigidity witness.
\end{theorem}

=======
We introduce the notion of a \emph{Terminal Rigidity Witness}: a physical system whose global
structure is maximally nontrivial, yet whose operational information capacity is identically zero.
We show that the classical Einstein--Rosen bridge (ERB), the maximal analytic extension of the
Schwarzschild solution, furnishes a canonical example of such a witness.
Formally, ERB induces a zero channel in the sense of transcript capacity, and thus constitutes
a terminal element in the preorder of admissible physical inference systems.
This provides a rigorous stress test for claims of nonlocal reconstruction, exotic communication,
and entanglement-based operational geometry.
\end{abstract}

\section{Background and Motivation}

Modern theoretical physics and computation are increasingly populated by claims of
\emph{exotic accessibility}: wormholes enabling communication, black hole interiors being
operationally reconstructible, or entanglement acting as a geometric bridge between distant
regions of spacetime.

At the same time, across logic, computer science, and thermodynamics, we encounter
persistent \emph{locality obstructions}: finite-variable logics cannot express global structure,
bounded transcripts cannot encode global state, and physical interfaces cannot exceed
entropy bounds.

The central question is therefore structural rather than technological:
\begin{quote}
Is there a physical system whose global structure is as extreme as possible,
yet whose operational information content is provably zero?
\end{quote}
We argue that the classical Einstein--Rosen bridge answers this question in the affirmative.

\section{Interfaces and Transcript Capacity}

We adopt an operational framework of finite-capacity inference.

\begin{definition}[Interface]
An interface mediates interaction between an unknown system state $X$ and an observer,
producing a transcript $Y_{1:T}$ via admissible observation rules.
\end{definition}

\begin{definition}[Transcript Capacity]
The transcript capacity over horizon $T$ is
\[
\TC(T) := \sup I(X; Y_{1:T}).
\]
\end{definition}

This quantity measures the maximal operational information extractable by any admissible observer.

\section{Physical Systems and the SIGC Preorder}

We model physical systems as triples.

\begin{definition}[Physical System]
A physical system is a triple
\[
S := (\mathcal{I}, \mathcal{L}, E)
\]
where $\mathcal{I}$ is a finite set of interfaces,
$\mathcal{L} = (Y_{1:T})$ is an interaction log,
and $E$ is total dissipated energy.
\end{definition}

Each system induces a channel $\mathcal{C}_S : X \to Y_{1:T}$.
Define $\SIGC$ as the preorder of channels modulo admissible simulation:
\[
\mathcal{C}_1 \preceq \mathcal{C}_2
\iff
\Capacity(\mathcal{C}_1) \le \Capacity(\mathcal{C}_2).
\]

\section{Terminal Rigidity Witnesses}

\begin{definition}[Terminal Rigidity Witness]
A system $S$ is a terminal rigidity witness if
\[
\Capacity(\mathcal{C}_S) = 0
\]
and $S$ possesses nontrivial global structure.
\end{definition}

Such systems sit at the minimal element of the SIGC preorder: no admissible refinement can extract
information from them.

\section{The Einstein--Rosen Bridge}

Let $(M,g)$ be the maximal analytic extension of the Schwarzschild spacetime.
This spacetime contains two asymptotically flat regions $\Sigma_-$ and $\Sigma_+$ connected
by an Einstein--Rosen bridge.

\begin{definition}[ERB Channel]
Define the ERB channel
\[
\mathcal{C}_{\mathrm{ERB}} : \Sigma_- \to \Sigma_+
\]
as the induced causal communication map.
\end{definition}

\begin{lemma}[Topological Censorship]
Under global hyperbolicity and the averaged null energy condition,
no causal curve connects $\Sigma_-$ to $\Sigma_+$.
\end{lemma}

\begin{theorem}[ERB Zero Capacity]
\[
\Capacity(\mathcal{C}_{\mathrm{ERB}}) = 0.
\]
\end{theorem}

\begin{proof}
By topological censorship, all admissible causal trajectories fail to connect the two boundaries.
Hence outputs are statistically independent of inputs, so
$I(X;Y)=0$ for all encodings.
\end{proof}

\begin{corollary}
The Einstein--Rosen bridge is a terminal rigidity witness.
\end{corollary}

\section{Intuition: Geometry Without Information}

The ERB exhibits a striking asymmetry:
\begin{itemize}[leftmargin=*]
\item \textbf{Geometric richness}: nontrivial global topology and maximal analytic extension.
\item \textbf{Operational emptiness}: zero information capacity.
\end{itemize}
In plain terms, ERB is a \emph{perfect bridge that carries nothing}.

\section{Implications for Exotic Physics}

Any claim of:
\begin{itemize}[leftmargin=*]
\item traversable wormholes,
\item entanglement-based communication,
\item black hole interior reconstruction,
\end{itemize}
must violate at least one of:
\begin{itemize}[leftmargin=*]
\item locality,
\item thermodynamic calibration,
\item transcript monotonicity.
\end{itemize}

\section{Relation to the Viotropic Wall}

The Viotropic Wall theorem states that when informational demand exceeds transcript capacity,
observables become operationally independent.
The ERB is the extremal case: demand is nonzero, capacity is exactly zero.

\section{Conclusion}

The Einstein--Rosen bridge furnishes a canonical physical invariant:
a system of maximal global structure and zero operational content.
It functions as a terminal rigidity witness for all theories of physical inference.

\begin{thebibliography}{99}

\bibitem{EinsteinRosen1935}
A. Einstein and N. Rosen,
\emph{The Particle Problem in the General Theory of Relativity},
Physical Review \textbf{48}, 73--77 (1935).

\bibitem{Kruskal1960}
M. D. Kruskal,
\emph{Maximal Extension of Schwarzschild Metric},
Physical Review \textbf{119}, 1743--1745 (1960).

\bibitem{FriedmanSchleichWitt1993}
J. L. Friedman, K. Schleich, and D. M. Witt,
\emph{Topological Censorship},
Physical Review Letters \textbf{71}, 1486--1489 (1993).

\bibitem{MorrisThorne1988}
M. Morris and K. Thorne,
\emph{Wormholes in Spacetime and Their Use for Interstellar Travel},
American Journal of Physics \textbf{56}, 395--412 (1988).

\bibitem{Visser1996}
M. Visser,
\emph{Lorentzian Wormholes: From Einstein to Hawking},
Springer (1996).

\bibitem{Landauer1961}
R. Landauer,
\emph{Irreversibility and Heat Generation in the Computing Process},
IBM Journal of Research and Development \textbf{5}, 183--191 (1961).

\bibitem{Libkin2004}
L. Libkin,
\emph{Elements of Finite Model Theory},
Springer (2004).

\bibitem{Grohe2017}
M. Grohe,
\emph{Descriptive Complexity, Canonisation, and Definable Graph Structure Theory},
Cambridge University Press (2017).

\bibitem{Vasquez2026Viotropic}
I. F. Vasquez,
\emph{The Viotropic Wall: Finite Capacity, Thermodynamic Calibration,
and the Limits of Local Inference},
2026. Manuscript.

\end{thebibliography}

>>>>>>> f178752 (docs: restore canonical ERB LaTeX source)
\end{document}

